% ============================================================================
% EXERCISE C: SOWING DATE FACTORIAL
% ============================================================================
% Duration: 90 minutes (hands-on; Module 5 support)
% Learning Goals: Phenology effects; sowing date optimization; yield/quality trade-offs
% Prerequisites: Module 3, Module 5

\chapter{Exercise C: Sowing Date Factorial}
\label{ch:exercise-c}

\section{Overview}

\textbf{Objective:} Run three near-solved simulations with different sowing date windows and compare the effects on crop development, yield, and grain quality.

\begin{center}
\begin{tabular}{l|l|l|l}
\textbf{} & \textbf{Early May} & \textbf{Mid June} & \textbf{Late July} \\
\hline
File & \filename{Exercise\_C\_EarlyMay.apsimx} & \filename{Exercise\_C\_MidJun.apsimx} & \filename{Exercise\_C\_LateJul.apsimx} \\
Sowing window & 1 May -- 31 May & 15 Jun -- 15 Jul & 1 Aug -- 31 Aug \\
Duration & 90 minutes & & \\
\end{tabular}
\end{center}

\textbf{Simulated season:} 1949-07-01 to 1950-12-31 (one full growing season)

\textbf{Location:} Northam, Western Australia

\textbf{Deliverables:} Summary table, phenology comparison chart, yield table, written analysis

\section{Background}

Sowing date is one of the most powerful management levers available to a wheat farmer. Because wheat development is driven by thermal time (GDD), shifting the sowing date by 4--6 weeks moves the entire crop calendar, including the most stress-sensitive stages (anthesis, grain fill), into different temperature and rainfall conditions.

\begin{keyconceptbox}
\textbf{The core question:} For a given location and season, which sowing date best aligns the critical period (booting to grain fill) with favorable temperatures and soil water availability?
\end{keyconceptbox}

\section{Task 1: Open and Review the Simulation Files}

\begin{enumerate}
  \item Open \filename{Exercise\_C\_EarlyMay.apsimx} in APSIM
  \item Click the \apsimterm{Manager} component in the simulation tree
  \item Locate the sowing rule. You will see two parameters:
    \begin{itemize}
      \item \variable{StartDate}: ``1-may''
      \item \variable{EndDate}: ``31-may''
    \end{itemize}
  \item Note the sowing conditions: ESW $>$ 100 mm AND 7-day rainfall $>$ 25 mm
  \item Open \filename{Exercise\_C\_MidJun.apsimx} and \filename{Exercise\_C\_LateJul.apsimx}; verify that all other parameters (soil, variety, N rate, weather file) are identical --- only the sowing window differs
\end{enumerate}

% ---- IMAGE PLACEHOLDER ----
\begin{figure}[H]
  \centering
  \fbox{\parbox{0.85\textwidth}{\centering\vspace{1.5cm}
    \textit{[IMAGE: APSIM screenshot showing the Manager component properties panel for one of the Exercise C files.}\\
    \textit{Annotate: StartDate and EndDate parameters highlighted in the sowing rule.]}
  \vspace{1.5cm}}}
  \caption{Manager component in Exercise C showing the sowing window parameters. Only StartDate and EndDate differ between the three files.}
  \label{fig:exercise-c-manager}
\end{figure}
% ---- END PLACEHOLDER ----

\section{Task 2: Run All Three Simulations and Export Results}

\begin{enumerate}
  \item Open \filename{Exercise\_C\_EarlyMay.apsimx}; press \textbf{Run} (F5)
  \item Wait for completion (progress bar at bottom)
  \item Export results: Results tab $\to$ right-click $\to$ Export $\to$ Excel
  \item Save as \filename{C\_EarlyMay\_Results.xlsx}
  \item Repeat for \filename{Exercise\_C\_MidJun.apsimx} and \filename{Exercise\_C\_LateJul.apsimx}
\end{enumerate}

\begin{tipbox}
\textbf{Efficiency tip:} You can run all three simulations in sequence before exporting. APSIM stores results from each run separately in its DataStore.
\end{tipbox}

\section{Task 3: Phenological Comparison}

\subsection{Extract Key Dates}

For each simulation, find the date when each Zadok stage is first reached. In your exported Excel file:

\begin{enumerate}
  \item Add a helper column: \texttt{=YEAR(A2)} and \texttt{=MONTH(A2)} for date filtering
  \item Use \texttt{MINIFS} to find the date when \variable{[Wheat].Stage} first crosses each threshold:
    \begin{itemize}
      \item Emergence: Stage $\geq$ 10 (\texttt{=MINIFS(\$A:\$A, \$C:\$C, ">=10")} where column C = Stage)
      \item Booting: Stage $\geq$ 45
      \item Anthesis: Stage $\geq$ 60
      \item Physiological maturity: Stage $\geq$ 90
    \end{itemize}
\end{enumerate}

\subsection{Thermal Time Accumulation}

For each scenario, calculate cumulative GDD from sowing date to anthesis:
\begin{enumerate}
  \item Identify the sowing date (first row where Stage $>$ 0)
  \item Identify the anthesis date (first row where Stage $\geq$ 60)
  \item In a Python or Excel formula, sum the daily GDD column for that date range
\end{enumerate}

\begin{keyconceptbox}
\textbf{Expected result:} All three scenarios should accumulate similar GDD from sowing to anthesis ($\sim$850--950 GDD). The differences will be in \emph{how many calendar days} this takes --- cool early-winter sowing takes more calendar days to accumulate the same GDD as a warmer late-winter sowing.
\end{keyconceptbox}

\section{Task 3: Phenological Comparison}

\subsection{Key Dates}

Track dates for each sowing scenario:
\begin{itemize}
  \item Sowing date
  \item Emergence date (Zadok stage 10)
  \item Booting date (Zadok stage 40--49)
  \item Anthesis date (Zadok stage 60--69)
  \item Physiological maturity / Harvest date (Zadok stage 90)
\end{itemize}

\subsection{Thermal Time Accumulation}

Calculate cumulative GDD from sowing to anthesis. Compare across sowing dates.

\section{Task 4: Yield and Quality Comparison}

\subsection{Harvest Results Table}

For each scenario, extract from the row where Zadok Stage = 90 (harvest):

\begin{center}
\begin{tabular}{l|r|r|r}
\textbf{Metric} & \textbf{Early May} & \textbf{Mid June} & \textbf{Late July} \\
\hline
Actual sowing date & \textit{record} & \textit{record} & \textit{record} \\
Emergence date (Stage 10) & \textit{record} & \textit{record} & \textit{record} \\
Anthesis date (Stage 60) & \textit{record} & \textit{record} & \textit{record} \\
Maturity date (Stage 90) & \textit{record} & \textit{record} & \textit{record} \\
Season length (days) & \textit{calculate} & \textit{calculate} & \textit{calculate} \\
GDD sowing to anthesis & \textit{calculate} & \textit{calculate} & \textit{calculate} \\
Grain yield (kg/ha) & \textit{record} & \textit{record} & \textit{record} \\
Biomass at harvest (kg/ha) & \textit{record} & \textit{record} & \textit{record} \\
Harvest Index & \textit{calculate} & \textit{calculate} & \textit{calculate} \\
Grain protein (\%) & \textit{record (if available)} & & \\
\end{tabular}
\end{center}

\subsection{Expected Sowing Date Effects}

\begin{table}[H]
\centering
\caption{Expected outcomes by sowing date (indicative)}
\label{tab:sowing-date}
\begin{tabularx}{\textwidth}{l|X|X|X}
\toprule
\textbf{Metric} & \textbf{Early May} & \textbf{Mid June} & \textbf{Late July} \\
\midrule
Season length & Long ($\sim$230+ days) & Moderate ($\sim$195 days) & Short ($\sim$165 days) \\
Calendar days, sowing to anthesis & Many (cool autumn/winter) & Moderate & Few (warm spring) \\
Anthesis temperature & Cool (Aug) & Mild (Sep) & Warm (Oct--Nov) \\
Frost risk at booting & Higher (July booting) & Lower & Very low (Sep booting) \\
Heat stress at anthesis & Very low & Low & Moderate \\
Grain fill duration & Long & Moderate & Short \\
Grain yield & Moderate--high (frost risk) & Typically highest & Lower (short fill) \\
Grain protein & Low--moderate & Moderate & Higher (heat stress) \\
Harvest Index & Moderate--high & High ($\sim$0.45) & Lower ($\sim$0.38) \\
\bottomrule
\end{tabularx}
\end{table}

\section{Task 5: Analysis Questions}

Answer the following questions using your harvest results table and plots. Support each answer with specific numbers from your simulation output.

\begin{enumerate}
  \item \textbf{Calendar days vs.~thermal time:} Compare the number of calendar days from sowing to anthesis across the three scenarios. Then compare cumulative GDD from sowing to anthesis. Which is more similar across scenarios --- calendar days or GDD? What does this tell you about how APSIM models crop development?

  \item \textbf{Yield response:} Which sowing date produced the highest grain yield in this season? Using the ESW output and T$_{max}$ output, explain \emph{why} that scenario outperformed the others. Was it related to water availability, temperature, or both?

  \item \textbf{Critical period and stress:} For the Late July sowing, identify the date of anthesis in your output. What was the T$_{max}$ on that date? Was ESW above the 100 mm sowing threshold at that time? Does the output suggest water stress or heat stress contributed to any yield reduction?

  \item \textbf{Season length and grain fill:} Calculate season length (sowing to maturity) for each scenario. How does grain fill duration (maturity minus anthesis) change with sowing date? How might a shorter grain fill period explain differences in yield and grain size?

  \item \textbf{Harvest Index:} Calculate Harvest Index for each scenario (Grain mass / Biomass). Which scenario has the highest HI? What does a lower HI indicate about where assimilate ended up (straw vs.~grain)?

  \item \textbf{Practical recommendation:} Based on your simulation results for this one season (1949--1950), which sowing date would you recommend to a Northam wheat farmer in the WA Wheatbelt? Acknowledge at least one limitation of making this recommendation from a single-season analysis.
\end{enumerate}

\section{Task 6: Visualisation}

Create the following two plots:

\subsection{Crop Development Comparison (Line Plot)}

\begin{enumerate}
  \item In Excel (or Python), plot Zadok Stage (y-axis) vs.~Date (x-axis) for all three scenarios on the same axes
  \item Use different line colours/styles for each scenario
  \item Add vertical dashed lines at: frost risk window (Jul 1 -- Jul 31) and heat risk window (Oct 1 -- Nov 30)
  \item Add horizontal dashed lines at Stage 45 (Booting), Stage 60 (Anthesis), Stage 90 (Maturity)
  \item Title: ``Phenological Development by Sowing Date''
\end{enumerate}

% ---- IMAGE PLACEHOLDER ----
\begin{figure}[H]
  \centering
  \fbox{\parbox{0.85\textwidth}{\centering\vspace{1.5cm}
    \textit{[IMAGE: Student-generated Zadok stage vs.~date line plot for all three sowing date scenarios.}\\
    \textit{Annotate: frost risk window shaded in blue; heat risk window shaded in red.]}
  \vspace{1.5cm}}}
  \caption{Example student output: Zadok stage trajectory for three sowing dates. The frost and heat risk windows highlight the trade-off between early and late sowing.}
  \label{fig:exercise-c-phenology-plot}
\end{figure}
% ---- END PLACEHOLDER ----

\subsection{Yield and Harvest Index Comparison (Bar Chart)}

Create a grouped bar chart with:
\begin{itemize}
  \item X-axis: three sowing date scenarios
  \item Two bar groups: Grain yield (left axis, kg/ha) and Harvest Index (right axis, 0--1)
  \item Title: ``Grain Yield and Harvest Index by Sowing Date''
\end{itemize}

\section{Deliverable}

Submit a single Excel workbook (or Python notebook + PDF) containing:

\begin{enumerate}
  \item \textbf{Harvest results table} (Task 4): all metrics for three scenarios
  \item \textbf{Phenology comparison chart} (Task 6, Plot 1): Zadok stage vs.~date
  \item \textbf{Yield/HI bar chart} (Task 6, Plot 2): grouped bar chart
  \item \textbf{Written analysis} (Task 5): six questions, approximately 2--4 sentences each
  \item \textbf{Recommendation} (Task 5, Q6): one paragraph with supporting evidence
\end{enumerate}

\textbf{Total time:} 90 minutes

\textbf{Assessment criteria:}
\begin{itemize}
  \item Harvest table complete and values are plausible (20\%)
  \item Plots are correctly labelled with units and titles (20\%)
  \item Analysis questions answered with specific data values (40\%)
  \item Recommendation is evidence-based and acknowledges limitations (20\%)
\end{itemize}

% ============================================================================
% END EXERCISE C
% ============================================================================
